%%%%%%%%%%%%%%%%%%%%%%%%%%%%%%%%%%%%%%%%%%%%%%%%%%%%%%%%%%%%%%%%%%%%%%%%%%%%%%
%  qb-unit1.tex -- Unit I: Introduction to Digital Marketing; Content Marketing
%%%%%%%%%%%%%%%%%%%%%%%%%%%%%%%%%%%%%%%%%%%%%%%%%%%%%%%%%%%%%%%%%%%%%%%%%%%%%%
\chapter{Introduction to Digital Marketing and Content Marketing}

\textit{Syllabus:} Principles of Digital Marketing; Digital Marketing Channels;
Tools to Create Buyer Persona; Competitor Research Tools; Website Analysis
Tools. Content Marketing Concepts and Strategies; Planning, Creating,
Distributing and Promoting Content; Optimising Website UX and Landing Pages;
Measuring Impact; Metrics and Performance; Using Content Research for
Opportunities.

\textit{In the examination:} Part~B Question~2 is compulsory and is always set
from this unit, so Unit~I cannot be skipped by internal choice. It also
supplies the largest share of Part~A, because the marketing fundamentals ---
the mix, segmentation and research vocabulary --- sit here.

\parta

\begin{qlist}

\qq{25-1.1}{Differentiate between push and pull marketing.}{CIE-I, 3 Apr 2025;
CIE-I, 16 Apr 2026; SEE, Jun/Jul 2025, Q1.1 --- 2 M}
\textbf{Push marketing} takes the product to the customer: the seller creates
demand and pushes the offering through the channel --- trade promotions,
distributor incentives, cold e-mails, SMS blasts, display and interruptive ads.
The buyer has not asked for it. It suits new or unknown products that nobody is
searching for yet.
\textbf{Pull marketing} draws the customer to the brand: demand is created so
that the buyer initiates the search --- SEO, content marketing, organic social,
referrals, word of mouth. In one line: push interrupts, pull attracts.
\scheme{1 M for push with an example, 1 M for pull with an example.}

\qq{25-1.2}{Mention any four characteristics of a good content.}{CIE-I, 3 Apr
2025; CIE-I, 16 Apr 2026; SEE, Jun/Jul 2025, Q1.2 --- 2 M}
\begin{apts}
\item \textbf{Relevant and useful} --- solves a real problem of the defined
      target audience.
\item \textbf{Original and accurate} --- researched, credible, correctly
      sourced, not duplicated.
\item \textbf{Clear and well-structured} --- readable, scannable headings,
      short paragraphs, plain language.
\item \textbf{Engaging and actionable} --- holds attention and closes with a
      clear call to action.
\end{apts}
Also acceptable: SEO-optimised, shareable, visually supported, consistent with
brand voice, current and up to date.
\scheme{Any four characteristics, half a mark each. Naming alone is enough at
two marks.}

\qq{24-1.1}{Comment on the importance of demographic segmentation.}{SEE,
Sept/Oct 2024, Q1.1 --- 2 M}
Demographic segmentation divides a market by measurable population variables
--- age, gender, income, education, occupation, family size, religion. It is
important because the data is easy and cheap to obtain, is objective and
verifiable, correlates strongly with both needs and purchasing power, and is
directly usable as targeting criteria on every digital advertising platform. It
therefore allows precise message personalisation and cost-efficient media
selection, and it is the base layer on which psychographic and behavioural
segmentation are built.
\scheme{1 M for what it is, 1 M for two or more reasons why it matters.}

\qq{24-1.2}{Identify the 4 P's of marketing.}{SEE, Sept/Oct 2024, Q1.2 --- 2 M}
\begin{apts}
\item \textbf{Product} --- the goods or service offered, its features, quality,
      branding and after-sales support.
\item \textbf{Price} --- the amount charged and the pricing strategy, discounts
      and payment terms.
\item \textbf{Place} --- the distribution channels and locations through which
      the customer obtains it.
\item \textbf{Promotion} --- the communication mix used to inform and persuade:
      advertising, sales promotion, personal selling, publicity, digital.
\end{apts}
\scheme{Half a mark each. Naming the four is sufficient; a phrase of
explanation each earns full marks safely.}

\qq{24-1.3}{Assume you are sitting in the dining area in casual attire during
your vacation. Your uncle, a marketing manager, gives you a pen and asks you to
sell it right at that moment. How would you respond?}{SEE, Sept/Oct 2024, Q1.3
--- 2 M}
I would not begin by describing the pen. I would first ask questions to
uncover a need --- ``When did you last have to sign something important, and
what did you sign it with?'' --- and then position the pen against the need
that answer reveals, whether that is reliability at a client meeting, the
impression a good instrument makes, or simply having one to hand when it
matters. I would close by asking for the sale directly. The point being
demonstrated is \textbf{need-based selling}: establish the need first, present
the product as the solution second.
\scheme{1 M for asking questions/establishing the need before pitching, 1 M for
linking the product's benefit to that need and closing.}

\qq{24-1.4}{Identify two types of research data.}{SEE, Sept/Oct 2024, Q1.4
--- 2 M}
\textbf{Primary data} --- collected first-hand by the researcher for the
problem currently in hand, through surveys, interviews, focus groups,
observation or experiments. It is specific and current but costly and slow.
\textbf{Secondary data} --- already collected by someone else for another
purpose, such as census reports, industry and trade publications, company
records or existing Google Analytics data. It is fast and cheap but may be
dated or not exactly fit the problem.
\scheme{1 M each, with one line of description.}

\qq{24-1.5}{Provide your own psychographic and demographic profile as a
marketing person would assess you.}{SEE, Sept/Oct 2024, Q1.5 --- 2 M}
\textbf{Demographic:} 20 years old, unmarried, undergraduate engineering
student, resident of Bengaluru, urban middle-income household, no independent
income beyond a monthly allowance of roughly \Rs 5{,}000.
\textbf{Psychographic:} technology-forward and an early adopter of apps;
value-conscious and strongly discount-driven; achievement-oriented with
placement and higher studies as goals; interests in gaming, fitness and
coding; heavy smartphone and social media user who reads reviews before every
purchase; brand-aware but price-sensitive; prefers online purchase with
cash-on-delivery or UPI.
\scheme{1 M for a demographic profile using at least three standard variables,
1 M for a psychographic profile covering lifestyle, interests, values or
attitudes.}

\qq{24-1.6}{Highlight one advantage and one limitation of psychographic
segmentation.}{SEE, Sept/Oct 2024, Q1.6 --- 2 M}
\textbf{Advantage:} it explains \emph{why} people buy. By segmenting on values,
lifestyle, personality, interests and attitudes it enables emotionally
resonant messaging and genuine brand differentiation, and so builds stronger
loyalty than demographics alone can.
\textbf{Limitation:} the data is subjective and difficult to obtain. It cannot
be observed directly, must be inferred from surveys or behavioural proxies, is
expensive to collect and hard to validate, and the resulting segments are hard
to size and reach precisely --- which makes media buying against them less
reliable than against demographics.
\scheme{1 M each.}

\qq{24-1.7}{You want to test an app you have developed which addresses some
needs of engineering college students. How can you use digital marketing to get
initial quantitative feedback?}{SEE, Sept/Oct 2024, Q1.7 --- 2 M}
Run narrowly targeted paid campaigns on Instagram, Facebook and LinkedIn
filtered to the 18--24 age band and to specific engineering colleges, driving
traffic to a landing page carrying an embedded Google Form. Quantitative
feedback then comes from the numbers the campaign itself produces ---
click-through rate, cost per install, sign-up conversion rate, form response
distributions, in-app event completion and Day-7 retention from Firebase or
GA4 --- supplemented by A/B tests of two value propositions and an in-app
rating prompt. Every one of these is a countable metric, which is what
distinguishes quantitative from qualitative feedback.
\scheme{1 M for naming a targeted digital channel and mechanism, 1 M for
naming the numeric metrics that constitute the quantitative feedback.}

\qq{24-1.8}{List any four online marketing tools that can be useful for
beginners in the field of marketing.}{SEE, Sept/Oct 2024, Q1.8 --- 2 M}
\begin{apts}
\item \textbf{Google Analytics} --- website traffic, behaviour and conversion
      measurement.
\item \textbf{Google Keyword Planner} or \textbf{Google Trends} --- keyword
      demand and topic research.
\item \textbf{Canva} --- creative and social media design without a designer.
\item \textbf{Mailchimp} --- e-mail list management, campaigns and automation.
\end{apts}
Also acceptable: Hootsuite or Buffer for scheduling, SEMrush or Ubersuggest for
SEO, Google Search Console, Meta Business Suite, Hotjar.
\scheme{Any four, half a mark each.}

\qq{c1a1}{Differentiate between marketplace and market space.}{CIE-I, 3 Apr
2025; CIE-I, 16 Apr 2026 --- 2 M}
\textbf{Marketplace} is a physical or definite location at which buyers and
sellers meet to exchange goods and services --- a shop, a mall, a showroom, a
mandi. \textbf{Market space} is the virtual or conceptual environment in which
commercial value is created and exchanged through digital networks --- Amazon,
an app, an online store --- with no physical location, no fixed hours and no
geographic limit. In short: marketplace is physical, market space is digital.
\scheme{1 M each.}

\qq{c1a2}{What is the mantra of marketing?}{CIE-I, 3 Apr 2025 --- 2 M}
\textbf{Create, Communicate and Deliver Value to a Target market at a Profit.}
It compresses the whole discipline into one line: value must first be created
in the offering, then communicated so the market knows of it, then delivered so
the promise is kept --- all aimed at a chosen target market and all done
profitably, since value delivered at a loss is not marketing but charity.
\scheme{1 M for the mantra itself, 1 M for a brief expansion of what it means.}

\qq{c1a3}{Define affiliate marketing and give an example.}{CIE-I, 3 Apr 2025
--- 2 M}
Affiliate marketing is a performance-based model in which a business rewards
external partners --- affiliates --- with a commission for each visitor, lead or
sale their own promotional efforts deliver through a tracked link. The
advertiser pays only for results, which makes the risk almost entirely the
affiliate's. \textbf{Example:} Amazon Associates, where a blogger or YouTuber
earns a percentage of every purchase made through their referral link.
\scheme{1 M for the definition, 1 M for a correct named example.}

\qq{c1a4}{What is the first step of content marketing?}{CIE-I, 16 Apr 2026 ---
2 M}
Defining the goals and identifying the target audience. Everything downstream
--- topic selection, format, tone, channel and the metric used to judge success
--- depends on knowing what the content is meant to achieve and who it is meant
to reach, so content created before this step is decided is content produced
without a standard to be judged against.
\scheme{1 M for goal setting and audience identification, 1 M for why it comes
first.}

\qq{c1a5}{Mention any two core concepts of marketing.}{CIE-I, 16 Apr 2026 ---
2 M}
\textbf{Needs, wants and demands} --- a need is a basic requirement, a want is
the culturally shaped form that need takes, and a demand is a want backed by
the ability to pay. \textbf{Value and satisfaction} --- the customer's judgement
of the benefits received against the total cost paid, which determines whether
they buy again. Also acceptable: exchange and transactions; markets;
segmentation and targeting; the offering and the brand.
\scheme{1 M each, with a phrase of explanation.}

\qq{c1a6}{List any two principles of digital marketing.}{SEE, Oct/Nov 2023 ---
2 M}
\textbf{Know your customers} --- build the strategy around a researched buyer
persona rather than around the product. \textbf{Know how you are doing} ---
measure continuously, because the defining property of digital marketing is
that almost everything is trackable and therefore improvable. Also acceptable
from the five pillars: know your business; know the competition; know what you
want to achieve. Or, from the principle set: measurability, interactivity,
personalisation, permission-based targeting.
\scheme{1 M each.}

\qq{c1a7}{Differentiate between market segmentation and marketing mix.}{SEE,
Oct/Nov 2023 --- 2 M}
\textbf{Market segmentation} is the process of dividing a large heterogeneous
market into smaller homogeneous groups of buyers with similar needs, using
geographic, demographic, psychographic and behavioural bases, so that each
group can be served distinctly. \textbf{Marketing mix} is the set of
controllable tactical tools --- Product, Price, Place and Promotion --- blended
to produce the desired response from a chosen segment. In one line:
segmentation decides \emph{whom} to serve; the marketing mix decides
\emph{how} to serve them.
\scheme{1 M each. The ``whom against how'' contrast secures the second mark.}

\end{qlist}

\partb

\begin{qlist}

\qq{24-2a}{Explain any four key principles of digital marketing.}{SEE,
Sept/Oct 2024, Q2a --- 8 M}
\begin{apts}
\item \textbf{Customer-centricity.} Digital marketing is built outward from a
      documented buyer persona --- the customer's needs, context, objections
      and journey --- rather than inward from product features. Every channel,
      message and offer is chosen because it fits where that person already is.
\item \textbf{Measurability and data-driven decisions.} Almost every
      interaction is trackable, so decisions are made on analytics and KPIs
      rather than opinion. Budget follows evidence: a channel that does not
      show a return in the data is cut.
\item \textbf{Segmentation and personalisation.} The right message reaches the
      right person at the right moment, using demographic, behavioural and
      lifecycle data. Relevance, not volume, is what earns response.
\item \textbf{Value-first content.} Attention is earned by being useful ---
      solving a problem, answering a question, entertaining --- rather than
      bought by interruption. Content built this way keeps working long after
      the campaign ends.
\item \textbf{Integration and consistency across channels.} Search, social,
      e-mail, mobile and the website carry one brand voice and one promise, so
      the customer experiences a single journey rather than disconnected
      campaigns.
\item \textbf{Continuous testing and optimisation.} A/B testing, iteration and
      rapid feedback loops mean nothing is final; the campaign that runs in
      week eight is measurably better than the one launched in week one.
\item \textbf{Two-way engagement.} Digital media is a conversation ---
      comments, reviews, DMs and communities --- so listening and responding
      are part of the marketing, not an afterthought handled by support.
\item \textbf{Cost-efficiency and scalability.} Precise targeting removes
      wasted impressions, and a campaign that works can be scaled up or
      geographically extended in hours rather than quarters.
\end{apts}
\scheme{Any four principles, 2 M each: 1 M for correctly naming the principle,
1 M for explaining it. Do not write eight thin points --- the question asks for
four, explained.}

\qq{24-2b}{Explain any four digital marketing channels to spread awareness
about a new product.}{SEE, Sept/Oct 2024, Q2b --- 8 M}
\begin{apts}
\item \textbf{Search engine marketing.} Paid search captures demand
      immediately --- bidding on category and competitor keywords puts the new
      product in front of people already looking for that solution --- while
      SEO builds the durable organic visibility that carries awareness after
      the launch budget stops.
\item \textbf{Social media marketing.} Organic posts, Reels and Stories on
      Instagram, Facebook, YouTube and LinkedIn, amplified by paid reach and
      video-view campaigns, deliver mass awareness cheaply and allow precise
      interest and lookalike targeting for a product with no existing customer
      base.
\item \textbf{Content marketing.} Blog articles, explainer and demonstration
      videos, infographics, webinars and buying guides teach the market what
      the product is for. For a genuinely new category this is essential ---
      people cannot search for something they have no name for.
\item \textbf{E-mail marketing.} An announcement sequence to the existing
      opt-in list is the cheapest awareness channel available and reaches the
      warmest audience: existing customers who already trust the brand and are
      most likely to try and to refer.
\item \textbf{Influencer marketing.} Creators with an established, trusting
      niche audience lend credibility a new brand has not yet earned, and
      produce demonstration content at a fraction of studio cost.
\item \textbf{Display, video and programmatic advertising.} Banner, YouTube
      pre-roll and OTT placements build broad reach and visual recall, and
      support retargeting of everyone who visited the launch page.
\item \textbf{Affiliate and marketplace marketing.} Partners and marketplace
      listings promote on a pay-per-result basis, giving reach with almost no
      upfront risk.
\item \textbf{Mobile marketing.} App push notifications, SMS and WhatsApp
      broadcasts, in-app advertising and QR codes on packaging reach the
      customer on the device they actually use.
\end{apts}
\scheme{Any four channels, 2 M each: 1 M for naming and describing the channel,
1 M for showing specifically how it creates awareness for a \emph{new} product.}

\qq{25-2a}{Discuss how website analysis contributes to enhanced customer
experience and conversion optimization.}{SEE, Jun/Jul 2025, Q2a --- 8 M}
\textbf{Website analysis} is the systematic evaluation of a site's traffic,
technical health, content performance, usability and conversion paths, using
tools such as GA4, Google Search Console, PageSpeed Insights, Hotjar or
Microsoft Clarity, and SEMrush Site Audit.

\textit{Contribution to customer experience}
\begin{apts}
\item \textbf{It locates friction.} Exit-rate and funnel reports name the exact
      pages where visitors give up, turning a vague complaint that ``the site
      is confusing'' into a specific page to fix.
\item \textbf{It shows behaviour, not just numbers.} Heatmaps and session
      recordings expose dead clicks, rage clicks, ignored navigation and
      abandoned form fields --- the causes behind the exit statistics.
\item \textbf{It fixes performance.} Page-speed and Core Web Vitals analysis
      identifies slow-loading pages, the single most common destroyer of
      experience on mobile data connections.
\item \textbf{It enforces device reality.} Device, browser and screen-size
      reports reveal that most traffic is mobile and force responsive,
      mobile-first correction rather than desktop-led design.
\item \textbf{It surfaces unmet needs.} On-site search term reports say, in the
      customer's own words, what they came for and could not find --- a direct
      content and navigation backlog.
\end{apts}

\textit{Contribution to conversion optimisation}
\begin{apts}\setcounter{enumi}{5}
\item \textbf{It quantifies leakage.} Funnel visualisation gives the drop-off
      percentage at each step --- product view, add to cart, checkout, payment
      --- so effort goes where the loss is largest, not where opinion says.
\item \textbf{It segments quality of traffic.} Conversion rate by source,
      campaign, device and geography shows which channels bring buyers rather
      than visitors, redirecting budget to what converts.
\item \textbf{It enables controlled testing.} Analysis produces hypotheses that
      A/B and multivariate tests then prove or disprove --- headline, CTA
      wording and placement, form length, layout, imagery, pricing display.
\item \textbf{It closes the loop.} Conversion rate, average order value, cost
      per acquisition and revenue per visitor are re-measured after each
      change, so improvement is demonstrated rather than assumed.
\end{apts}

\textbf{Conclusion.} Website analysis replaces opinion-driven design with
evidence-driven design. The same body of data serves both goals at once:
removing friction raises satisfaction, and satisfaction is what a conversion
actually is --- so lower bounce, higher engagement, higher conversion rate and
higher lifetime value are outcomes of one activity, not two.
\scheme{2 M for defining website analysis and naming the tools, 3 M for the
customer-experience contributions, 3 M for the conversion contributions. A
concluding line linking the two earns the benefit of the doubt on the last
mark.}

\qq{25-2b}{Explain how content research tools help marketers discover new
opportunities and gain a competitive edge.}{SEE, Jun/Jul 2025, Q2b --- 8 M}
\textbf{Content research tools} are the software used to find out what an
audience is searching for, what competitors already rank for, and what is
gaining traction --- BuzzSumo, SEMrush and Ahrefs content-gap reports, Google
Trends, Google Keyword Planner, AnswerThePublic, Google Search Console, and
social listening platforms such as Brandwatch or Sprout Social.

\begin{apts}
\item \textbf{Topic and demand discovery.} Trend and volume data show which
      subjects are rising, which are seasonal and which are already saturated,
      so the brand publishes on an emerging topic \emph{before} the market
      crowds into it. Being early is itself the advantage.
\item \textbf{Keyword and intent research.} The tools reveal long-tail,
      low-difficulty keywords that carry real search volume and clear buying
      intent --- the queries a smaller brand can realistically win --- and
      allow each keyword to be mapped to a funnel stage.
\item \textbf{Content-gap analysis.} Comparing the brand's ranking keywords
      against three or four competitors produces an explicit list of topics
      the competitors earn traffic from and the brand does not. That list is a
      prioritised editorial backlog, generated from data rather than
      brainstorming.
\item \textbf{Competitor benchmarking.} The tools show which of a competitor's
      pages attract the most traffic, shares and backlinks, at what publishing
      cadence and in which formats --- so their proven winners inform strategy
      instead of guesswork.
\item \textbf{Audience question mining.} AnswerThePublic, Search Console query
      reports, on-site search logs, Reddit and Quora give the exact phrasing
      customers use for their problems, which makes copy resonate and captures
      voice and conversational search.
\item \textbf{Format and channel insight.} Engagement data shows whether a
      topic performs as short video, long-form guide, listicle or infographic,
      and on which platform --- preventing effort spent producing the right
      subject in the wrong form.
\item \textbf{Distribution and influencer discovery.} The tools identify who
      amplifies content in the niche and which publications link to it, so
      promotion is planned alongside creation rather than after it.
\item \textbf{Performance measurement and refresh.} Declining pages are flagged
      for updating, keyword cannibalisation between two of the brand's own
      pages is detected, and the next cycle's plan is corrected by the last
      cycle's results.
\end{apts}

\textbf{Competitive edge.} Taken together these give faster time-to-market on
emerging topics, better return on every keyword targeted, far less wasted
production effort, an editorial calendar that can be defended with evidence,
and an organic visibility position that competitors cannot buy their way past
quickly.
\scheme{1 M for naming the tools, 6 M for six distinct ways they create
opportunity (1 M each), 1 M for an explicit statement of the resulting
competitive edge.}

\qq{c1b1}{What is digital marketing? Explain the differences between traditional
marketing and digital marketing. Highlight any five key differences.}{SEE,
Oct/Nov 2023, Q2a --- 8 M}
\textbf{Digital marketing} is the promotion of products, services and brands
through digital channels and electronic devices --- search engines, websites,
social media, e-mail, mobile applications and display networks --- using data
to target, personalise and measure every interaction.

\begin{apts}
\item \textbf{Direction of communication.} Traditional marketing is one-way
      broadcast: the brand speaks and the audience listens. Digital marketing is
      two-way: comments, reviews, DMs and shares make the audience a
      participant. \emph{Example:} a newspaper advertisement cannot be replied
      to; an Instagram post is answered within minutes.
\item \textbf{Targeting precision.} Traditional targets a broad demographic
      inferred from the medium --- everyone who reads that paper. Digital
      targets by age, location, interest, behaviour, past purchase and lookalike
      modelling. \emph{Example:} a hoarding is seen by every passer-by; a Meta
      campaign reaches only women aged 25--34 in three postcodes.
\item \textbf{Measurability.} Traditional results are estimated through surveys
      and circulation figures, usually after the campaign. Digital results ---
      impressions, clicks, conversions, revenue, ROAS --- are visible in real
      time and attributable to the individual advertisement.
\item \textbf{Cost and accessibility.} Traditional demands large fixed
      expenditure that excludes small firms. Digital can start at a few hundred
      rupees a day and scale with results, which puts a startup and a
      multinational in the same auction.
\item \textbf{Speed and flexibility of change.} A printed or televised campaign
      is fixed once released; a digital campaign's budget, creative, audience
      and landing page can be changed mid-flight in response to the data.
\item \textbf{Reach.} Traditional reach is geographically bounded by the medium;
      digital reach is global and continuous, available at any hour.
\item \textbf{Personalisation.} Traditional delivers one identical message to
      everyone. Digital delivers dynamically different content, offers and
      recommendations to each segment or individual.
\end{apts}
\scheme{2 M for defining digital marketing, 5 M for five clear differences at
1 M each with an example, 1 M for presentation --- a comparison table earns this
easily. The question says ``any five'', so five differences properly explained
beat seven listed thinly.}

\qq{c1b2}{Describe the key principles of digital marketing and explain how they
differ from traditional marketing approaches. Provide examples to illustrate
your points.}{CIE-I, 1 Jul 2024 --- 10 M}
\textit{The principles} --- the five pillars first, then the operating
principles.
\begin{apts}
\item \textbf{Know your business.} Be clear about the offering, the value
      proposition and the commercial objective before choosing a channel. Where
      traditional marketing began with a media plan, digital begins with a
      measurement plan.
\item \textbf{Know the competition.} Competitive research is continuous and
      largely public --- their keywords, ads, content and backlinks are all
      observable through SEMrush or the Meta Ad Library, which had no
      traditional equivalent.
\item \textbf{Know your customers.} Strategy is built on a researched buyer
      persona drawn from analytics and social insight, not on the broad
      demographic that a newspaper's readership implied.
\item \textbf{Know what you want to achieve.} Objectives are set as SMART,
      channel-level KPIs --- cost per lead, ROAS --- rather than as an aspiration
      such as ``increase brand awareness''.
\item \textbf{Know how you are doing.} Continuous measurement and correction
      replace the post-campaign review. \emph{Example:} a Google Ads campaign
      whose budget shifts daily towards the better-converting keyword.
\item \textbf{Measurability.} Every impression, click and conversion is
      countable and attributable, so spend is defended with evidence.
      \emph{Traditional contrast:} a television spot's effect is inferred from
      a brand-lift survey weeks later.
\item \textbf{Interactivity and two-way engagement.} The audience replies,
      reviews and creates content. \emph{Example:} a brand answering complaints
      on X within the hour --- impossible in print.
\item \textbf{Personalisation.} Content adapts to the individual.
      \emph{Example:} Netflix's per-user home screen and Amazon's ``customers
      also bought'', against one identical billboard for all.
\item \textbf{Permission-based targeting.} Digital largely operates on consent
      --- opt-in lists, notification permissions, cookie consent --- where
      traditional simply interrupted.
\item \textbf{Iteration and agility.} Continuous A/B testing means the campaign
      improves while it runs; a printed campaign cannot be edited after it
      ships.
\item \textbf{Cost-efficiency and scalability.} Low entry cost, precise
      targeting and instant scaling replace large fixed media buys.
\end{apts}
\scheme{10 M. Roughly 5 M for the principles named and explained and 5 M for
the contrast with traditional marketing plus examples. The question has three
demands --- describe, differentiate, exemplify --- and marks are distributed
across all three, so an answer with no examples caps at about two-thirds.}

\qq{c1b3}{Identify and explain the various entities that can be marketed,
providing relevant examples for each category.}{CIE-I, 3 Apr 2025 --- 10 M}
Marketing applies to ten classes of entity, not merely to physical products.
\begin{apts}
\item \textbf{Goods} --- tangible physical products, the largest share of most
      economies' production and marketing effort. \emph{Example:} cars, mobile
      phones, packaged food.
\item \textbf{Services} --- intangible, perishable, inseparable offerings
      produced and consumed simultaneously. \emph{Example:} airlines, banking,
      hotels, salons, software as a service.
\item \textbf{Experiences} --- staged and orchestrated combinations of goods and
      services sold for the experience itself. \emph{Example:} Disneyland, an
      adventure trek, a music festival.
\item \textbf{Events} --- time-bound occasions marketed to attendees, sponsors
      and broadcasters. \emph{Example:} the Olympics, the IPL, a trade fair.
\item \textbf{Persons} --- celebrity, professional and personal branding.
      \emph{Example:} sportspersons, actors, musicians and consultants building
      a marketable public identity.
\item \textbf{Places} --- cities, states and countries marketed for tourism,
      residence and investment. \emph{Example:} ``Incredible India'', Dubai's
      tourism campaigns, state investor summits.
\item \textbf{Properties} --- intangible rights of ownership in real estate or
      financial assets. \emph{Example:} apartments, land, shares, mutual funds.
\item \textbf{Organizations} --- institutions marketing their own image and
      reputation to build preference and attract talent. \emph{Example:} a
      university, a hospital, a corporate employer brand.
\item \textbf{Information} --- knowledge produced, packaged and sold.
      \emph{Example:} textbooks, market research reports, online courses,
      newspapers.
\item \textbf{Ideas} --- causes, behaviours and social propositions promoted for
      adoption, which is the domain of social marketing. \emph{Example:} ``Swachh
      Bharat'', anti-smoking and road-safety campaigns, organ donation.
\end{apts}
\scheme{10 M --- ten entities at 1 M each, every one requiring both the
explanation and a relevant example. Naming the ten without examples earns about
half, since the question explicitly asks for examples for each category.}

\qq{c1b4}{How have the 4Ps of marketing evolved in the digital era? Provide
examples of companies effectively leveraging digital strategies for each
element.}{CIE-I, 3 Apr 2025 --- 10 M}
\begin{apts}
\item \textbf{Product --- from standardised to digital, personalised and
      continuously updated.} Products are now often intangible, delivered as
      software and improved after purchase; mass customisation replaces mass
      production, and user data feeds the next version.
      \emph{Example:} Spotify, whose Discover Weekly playlist is a product
      generated individually for each listener from their own listening data;
      Tesla, which improves a car already sold through over-the-air updates.
\item \textbf{Price --- from fixed printed price lists to dynamic, transparent
      and personalised pricing.} Prices adjust algorithmically to demand,
      competitor movement, inventory and user behaviour, while comparison sites
      make the whole market's pricing visible to the buyer, compressing margins
      and forcing value-based rather than cost-plus pricing. Subscription and
      freemium models replace one-time purchase.
      \emph{Example:} Amazon, which reprices millions of items many times a day;
      Uber's surge pricing; Netflix's tiered subscription.
\item \textbf{Place --- from physical distribution to omnichannel and direct
      digital delivery.} The shelf is replaced by the search result, the app and
      the marketplace listing; intermediaries are disintermediated as brands sell
      direct; and the customer expects to move between channels without losing
      continuity.
      \emph{Example:} Nike's direct-to-consumer app and website with
      click-and-collect linking online and store; Zara's integrated inventory
      allowing online return in a physical shop.
\item \textbf{Promotion --- from interruptive broadcast to targeted, interactive
      and measurable engagement.} Promotion becomes content, community and
      conversation --- SEO, social, influencer, e-mail and personalised
      advertising --- all of it measurable and adjustable in flight.
      \emph{Example:} Coca-Cola's ``Share a Coke'', which converted packaging
      into user-generated social content; Zomato's social media voice, which
      functions as continuous low-cost promotion.
\end{apts}

\textit{The extended mix.} Digital conditions have also promoted three further
Ps to first-class status --- \textbf{People} (community managers, reviewers and
influencers who now speak for the brand), \textbf{Process} (frictionless
checkout, one-click ordering, returns) and \textbf{Physical evidence} (website
design, reviews, ratings and trust signals, which are the only ``premises'' an
online buyer ever sees).
\scheme{10 M. 2 M for each of the four Ps --- 1 M for the nature of the
evolution and 1 M for a named company example --- and 2 M for a broader
observation such as the extended 7 Ps or a concluding synthesis. Examples are
explicitly demanded and carry half the marks.}

\qq{c1b5}{A retail company using newspaper ads, billboards and TV commercials is
considering a shift to digital marketing. Compare the benefits and limitations
of both approaches and suggest a strategy for transitioning.}{CIE-I, 3 Apr 2025;
CIE-I, 16 Apr 2026 --- 10 M}
\textit{Traditional marketing --- benefits:} very wide mass reach through
television and press; high perceived credibility and trust, particularly among
older audiences; strong local impact from hoardings and regional press;
excellent for broad brand building; and no dependence on the audience owning a
device or a connection.

\textit{Traditional marketing --- limitations:} expensive, with a high fixed
entry cost; results are difficult to measure and attribute; targeting is coarse;
communication is one-way; production and booking lead times are long, so nothing
can be corrected once released.

\textit{Digital marketing --- benefits:} far lower entry cost and spend that
scales with results; precise targeting by demography, interest, behaviour and
location; complete measurability and attribution; two-way interaction and
community; global and round-the-clock reach; rapid A/B testing and mid-flight
correction; and rich remarketing to people who have already shown interest.

\textit{Digital marketing --- limitations:} requires technical skill and tooling;
intense competition and rising auction costs; privacy regulation and consent
constraints; dependence on internet access, which limits rural and older
segments; platform dependence and algorithm volatility; and exposure to public
negative feedback and ad fraud.

\textit{Transition strategy}
\begin{apts}
\item \textbf{Audit and set objectives.} Establish the current cost per customer
      from traditional spend as the baseline, and set SMART digital objectives
      against it.
\item \textbf{Build the digital foundation first.} A responsive website with
      analytics, a Google Business Profile, and social profiles --- there is no
      point advertising to a destination that does not exist.
\item \textbf{Instrument measurement before spending.} GA4, conversion tracking
      and a CRM, so the first rupee spent is already measurable.
\item \textbf{Research the audience and build personas} from existing customer
      data and analytics rather than assumption.
\item \textbf{Run a small pilot in parallel, not a switch.} Keep traditional
      running and divert 10--20\% of budget into two digital channels for a
      quarter, so the comparison is evidence rather than faith.
\item \textbf{Start with high-intent channels.} Search advertising and local SEO
      capture demand the existing brand already generates, which produces the
      quickest demonstrable return and secures internal support.
\item \textbf{Add content and social} to build the owned audience that lowers
      paid dependence over time.
\item \textbf{Train or hire, and choose tooling.} The commonest cause of failed
      transition is capability, not budget.
\item \textbf{Integrate the two rather than opposing them.} Put QR codes and
      short URLs on print and hoardings, run television and search together
      around the same campaign, and measure the offline-to-online lift.
\item \textbf{Review, reallocate and scale.} Compare cost per acquisition
      channel by channel each quarter and move budget progressively towards
      what pays, retaining traditional media only where it still outperforms.
\end{apts}
\scheme{10 M. Roughly 4 M for the comparison --- benefits and limitations of
both, best set out as a table --- and 6 M for a phased, sensible transition
strategy. An abrupt ``stop traditional and go digital'' recommendation loses
marks; the graded, measured, parallel-run approach is what the scheme rewards.}

\qq{c1b6}{Your company has a limited marketing budget and must choose between
SEO, Pay Per Click and Social Media Marketing to promote a new high-tech gadget
aimed at young professionals. Which channel(s) would you prioritise and why?
Justify with specific examples.}{CIE-I, 1 Jul 2024 --- 10 M}
\textit{Compare the three against the situation.}
\begin{apts}
\item \textbf{SEO.} Free per click and compounding, with high buyer intent, but
      slow --- three to six months to rank --- and dependent on existing search
      demand.
\item \textbf{PPC.} Instant visibility, precise control, and captures intent at
      the moment of search, but costs money per click continuously and stops
      dead when the budget stops.
\item \textbf{SMM.} Strong for visual demonstration, discovery and virality
      among a young audience, with precise interest targeting and a low cost per
      thousand impressions, but low immediate intent --- nobody opened Instagram
      to buy a gadget.
\end{apts}

\textit{The decisive argument.} A brand-new gadget has a demand problem before
it has a conversion problem: nobody is searching for a product category they do
not yet know exists. SEO harvests existing demand and PPC harvests existing
search intent --- so neither can carry a launch on its own. Demand must first be
\emph{created}, and that is what social does.

\textit{Recommendation --- a weighted split, not a single channel.}
\begin{apts}\setcounter{enumi}{3}
\item \textbf{Social media marketing --- about 50\% of budget.} Instagram and
      YouTube short-form demonstration video, precisely targeted at young
      professionals by interest and job title, plus micro-influencer
      unboxings, which supply the credibility a new brand lacks.
      \emph{Example:} boAt built an entire audio brand in India primarily
      through social and influencer marketing rather than search.
\item \textbf{PPC --- about 35\%.} Bid on category and competitor keywords to
      capture the buyer who is comparison-shopping --- ``best noise cancelling
      earbuds under 5000'' --- and, critically, run retargeting against everyone
      the social campaign reached, which is where most of the conversion will
      actually occur. \emph{Example:} bidding on a competitor's brand name to
      appear beside it.
\item \textbf{SEO --- about 15\% now, rising later.} Begin immediately because
      it is slow: publish comparison and buying-guide content, optimise the
      product pages, and build the organic asset that will reduce paid
      dependence in months six to twelve. \emph{Example:} a ``best earbuds for
      commuting'' guide that keeps earning traffic long after the launch
      campaign ends.
\item \textbf{The sequencing logic.} Social creates awareness \ra\ PPC and
      retargeting convert the aware \ra\ SEO progressively harvests the demand
      the first two created, lowering cost per acquisition over time.
\item \textbf{Measurement and reallocation.} Judge each channel on cost per
      acquisition and ROAS rather than on clicks, review fortnightly, and shift
      budget as the launch matures --- the correct split at month one is not the
      correct split at month nine.
\end{apts}
\scheme{10 M. Roughly 3 M for comparing the three channels on the right
criteria (speed, cost model, intent, longevity), 5 M for a justified
recommendation with a clear reason tied to \emph{this} product and audience,
and 2 M for specific examples. The answer must commit to a priority --- a
balanced description of all three with no decision loses the justification
marks.}

\qq{c1b7}{A newly launched e-commerce company wants to build brand awareness
through content marketing. Suggest a step-by-step strategy, including content
types and distribution channels.}{CIE-I, 3 Apr 2025 --- 10 M}
\begin{apts}
\item \textbf{Goal setting.} Define what awareness means numerically --- for
      example branded search volume, reach and new-visitor sessions --- and set
      SMART targets against a baseline.
\item \textbf{Audience mapping.} Build buyer personas from market research,
      competitor audiences and early analytics; map their questions to the
      awareness, consideration and decision stages.
\item \textbf{Content ideation and planning.} Keyword and topic research,
      content-gap analysis against competitors, and an editorial calendar fixing
      topic, format, owner, channel and publication date.
\item \textbf{Content creation.} Produce to a defined quality standard ---
      original, useful, accurate, on-brand --- with a house style guide so output
      is consistent across writers.
\item \textbf{Content types.} Blog articles and buying guides for search;
      short-form video and Reels for discovery; infographics for shareability;
      podcasts and webinars for depth and authority; e-books and checklists as
      gated lead magnets; and user-generated content and reviews for social
      proof.
\item \textbf{Distribution --- owned media.} The website and blog, the e-mail
      list, the app and the brand's own social channels: no media cost, full
      control, and the audience that compounds.
\item \textbf{Distribution --- earned media.} SEO rankings, press coverage,
      guest posts, community and forum participation, shares and word of mouth:
      the most credible and the least controllable.
\item \textbf{Distribution --- paid media.} Promoted posts, search and display
      advertising, and influencer partnerships used to give the best-performing
      organic content additional reach rather than to prop up weak content.
\item \textbf{Amplification and repurposing.} Turn one substantial asset into
      many --- a webinar becomes a blog post, five Reels, an infographic and an
      e-mail --- and re-promote proven performers instead of always creating new.
\item \textbf{Evaluation and improvement.} Track reach and impressions,
      engagement rate, bounce rate, time on page, backlinks earned, branded
      search lift, leads and conversion; review monthly, cut formats that do not
      perform, refresh decaying content and feed the findings into the next
      quarter's calendar.
\end{apts}
\scheme{10 M --- ten sequenced steps at 1 M each. The question names three
things that must appear explicitly: the step-by-step structure, the content
types, and the distribution channels grouped as owned, earned and paid. Missing
any one of the three costs its marks however good the rest is.}

\qq{c1b8}{A newly launched health and wellness startup offering organic lifestyle
products aims to build brand awareness through content marketing. Design a
step-by-step content marketing strategy, outlining content formats, content
themes, distribution platforms (owned, earned, paid) and metrics for
performance evaluation.}{CIE-I, 16 Apr 2026 --- 10 M}
This question asks the same method as the previous one but names the four
deliverables that must appear. Answer under those four headings after a short
goal-and-audience preamble.

\textbf{Step 1 --- Goals.} Increase branded search volume by 200\% and reach
100{,}000 relevant users in six months; build an e-mail list of 10{,}000
opted-in subscribers as the owned asset.

\textbf{Step 2 --- Audience.} Health-conscious urban adults aged 25--45,
early-stage wellness adopters and new parents, who distrust unsubstantiated
health claims and respond to evidence, transparency and visible sourcing.

\begin{apts}
\item \textbf{Content formats.} Long-form blog articles and buying guides for
      search visibility; short-form video and Reels showing routines, recipes
      and product use; infographics on nutrition and ingredient comparisons;
      podcasts and expert interviews with nutritionists and doctors for
      authority; downloadable e-books, meal plans and habit trackers as gated
      lead magnets; and customer transformation stories and user-generated
      content for social proof.
\item \textbf{Content themes.} Education --- what organic certification actually
      means, how to read an ingredient label; Lifestyle --- routines, recipes,
      sleep, stress and fitness integration; Trust and transparency --- farm and
      supply-chain stories, certifications, third-party test results;
      Problem-solving --- content mapped to specific concerns such as gut health
      or immunity; Community --- challenges, testimonials and expert Q\&A.
\item \textbf{Owned distribution.} Website and SEO-optimised blog; e-mail
      newsletter and nurture sequences; Instagram, YouTube and WhatsApp
      channels; packaging inserts carrying QR codes to recipe content.
\item \textbf{Earned distribution.} Organic search rankings; wellness and
      parenting community participation; PR in health publications; expert and
      dietitian endorsements; customer reviews and organic shares.
\item \textbf{Paid distribution.} Meta and YouTube promotion of the
      best-performing organic content; search advertising on high-intent
      commercial keywords; micro-influencer partnerships with genuine wellness
      practitioners rather than general lifestyle creators.
\item \textbf{Metrics --- reach and visibility.} Impressions, reach, new users,
      organic sessions, keyword rankings, branded search volume.
\item \textbf{Metrics --- engagement.} Engagement rate, average engagement time,
      scroll depth, video completion rate, shares and saves.
\item \textbf{Metrics --- conversion and value.} Lead-magnet downloads, e-mail
      sign-up rate, conversion rate, cost per lead, revenue attributed to
      content, and customer lifetime value by acquisition content.
\item \textbf{Governance and compliance.} Health claims must be substantiated
      and legally compliant; every claim needs a source, and an expert review
      step belongs in the workflow. This is what distinguishes a wellness
      content strategy from a generic one.
\item \textbf{Review cycle.} Monthly performance review against the metrics
      above, quarterly content-gap re-analysis and refresh of decaying pages,
      with budget shifted towards the formats and themes that convert.
\end{apts}
\scheme{10 M. The scheme follows the four named deliverables: roughly 2 M each
for formats, themes, distribution across owned/earned/paid, and metrics, plus
2 M for the step structure and goal/audience framing. Answer under the four
headings the question gives --- examiners mark against the words in the
question.}

\qq{c1b9}{Imagine you are a digital marketing strategist at GreenTech
Innovations, launching a line of eco-friendly kitchen appliances called
``EcoChef''. Design a comprehensive content marketing strategy, including
specific goals, types of content, distribution channels and methods of
engagement.}{CIE-I, 1 Jul 2024 --- 10 M}
\begin{apts}
\item \textbf{Goals.} 50{,}000 website visits and 5{,}000 e-mail subscribers in
      the first six months; 1{,}000 pre-orders at launch; and top-three
      rankings for ``energy efficient kitchen appliances'' within a year --- each
      SMART and each measurable in GA4.
\item \textbf{Audience.} Environmentally conscious urban homeowners aged 28--45
      with above-median income; secondary audiences are new homeowners
      furnishing a kitchen and interior designers who specify appliances. Their
      decisive question is whether ``eco'' costs more than it saves.
\item \textbf{Positioning and content angle.} Lead with quantified savings ---
      units of electricity and rupees saved per year --- rather than with virtue.
      The strongest content asset is therefore a \emph{savings calculator}, not
      a blog post.
\item \textbf{Content types --- educational.} Blog guides on cutting kitchen
      energy use, explainers on energy ratings, and comparison articles against
      conventional appliances.
\item \textbf{Content types --- demonstrative.} Product demonstration and
      installation videos, a YouTube series on sustainable cooking, and
      before-and-after electricity-bill case studies from real customers.
\item \textbf{Content types --- interactive and gated.} The energy-savings
      calculator, a downloadable ``sustainable kitchen'' e-book, and a kitchen
      carbon-footprint quiz --- all capturing e-mail addresses.
\item \textbf{Distribution --- owned.} Website and SEO blog, e-mail nurture
      sequences, YouTube channel, Instagram and Pinterest (both strongly visual
      and strongly kitchen-relevant).
\item \textbf{Distribution --- earned and paid.} Earned: organic search,
      sustainability and home-improvement press, green certifications, customer
      reviews. Paid: search advertising on high-intent appliance keywords,
      YouTube pre-roll on cooking content, and partnerships with sustainability
      and home-design creators.
\item \textbf{Methods of engagement.} An ``EcoChef Challenge'' inviting users to
      post their reduced bills with a branded hashtag; live cooking sessions
      with chefs; a referral programme; responsive community management; a
      loyalty scheme rewarding reviews and repeat purchase; and user-generated
      kitchen photographs featured on the brand's channels.
\item \textbf{Measurement and iteration.} Track traffic, keyword rankings,
      engagement rate, calculator completions, e-mail sign-ups, pre-orders, cost
      per acquisition and content-attributed revenue; review monthly; double
      down on the formats that convert and retire those that do not.
\end{apts}
\scheme{10 M. The question names four deliverables --- goals, content types,
distribution channels, engagement methods --- so roughly 2--3 M each with the
remainder for measurement. The answer must be specific to EcoChef; a generic
content strategy with the brand name inserted twice will not earn the
application marks.}

\qq{c1b10}{Identify and explain at least five key metrics used to measure the
impact and performance of content marketing. How can these metrics influence
future content strategies?}{CIE-I, 1 Jul 2024 --- 10 M}
\textit{Reach and visibility metrics}
\begin{apts}
\item \textbf{Traffic --- users, sessions and pageviews by piece.} Measures how
      many people the content reached and through which channel.
      \emph{Influence on strategy:} the topics and formats that consistently
      draw traffic define the next quarter's calendar; those that draw none are
      retired rather than repeated.
\item \textbf{Organic search performance --- impressions, rankings, click-through
      rate from the SERP.} Measures durable discoverability.
      \emph{Influence:} pages with high impressions but low CTR need better
      titles and meta descriptions, not more content --- a cheap, high-return
      correction that only this metric reveals.
\end{apts}

\textit{Resonance metrics}
\begin{apts}\setcounter{enumi}{2}
\item \textbf{Engagement rate --- shares, saves, comments, likes, average
      engagement time.} Measures whether the content mattered to those it
      reached. \emph{Influence:} high shares and saves identify the formats
      worth repurposing and worth promoting with paid budget.
\item \textbf{Scroll depth and video completion rate.} Measures whether the
      content held attention to the end. \emph{Influence:} consistent drop-off
      at 30\% indicates the piece is too long or front-loaded with preamble,
      which changes how the next piece is structured.
\end{apts}

\textit{Friction metrics}
\begin{apts}\setcounter{enumi}{4}
\item \textbf{Bounce rate and exit rate by page.} Measures mismatch between the
      expectation that brought the visitor and what the page delivered.
      \emph{Influence:} a high bounce rate on pages arriving from one particular
      channel means the content and the channel are mismatched, so either the
      targeting or the piece must change.
\end{apts}

\textit{Business-outcome metrics}
\begin{apts}\setcounter{enumi}{5}
\item \textbf{Lead generation --- downloads, sign-ups, enquiries per piece.}
      Measures whether content produces a contactable prospect.
      \emph{Influence:} it identifies which topics deserve a gated asset built
      around them.
\item \textbf{Conversion rate and assisted conversions.} Measures content's
      contribution to revenue, including pieces that never get last-click
      credit. \emph{Influence:} prevents the classic mistake of killing
      top-of-funnel content that assists many conversions but closes few.
\item \textbf{Content ROI --- revenue attributed against production and
      promotion cost.} Measures whether the programme pays.
      \emph{Influence:} it settles budget arguments and sets the ceiling on
      what a piece of content is worth producing.
\item \textbf{Backlinks and brand mentions earned.} Measures authority
      accumulated. \emph{Influence:} identifies which content formats ---
      usually original data and research --- are worth the disproportionate
      effort they cost.
\end{apts}

\textbf{How metrics shape strategy overall.} Together they convert content from
a creative activity into a managed portfolio: double down on proven topics and
formats, fix rather than replace underperforming pages, refresh decaying
content, reallocate promotion budget to what earns it, and stop producing what
the data shows nobody reads --- so each cycle is measurably better informed than
the last.
\scheme{10 M. ``At least five'' means five is the floor: roughly 5--6 M for the
metrics named and explained and 4--5 M for the second half --- how each reading
changes the next content decision. The second half is the part candidates
routinely omit, and it carries close to half the marks.}

\qq{c1b11}{List and explain various pricing strategies used in marketing
management with suitable examples.}{SEE, Oct/Nov 2023, Q2b --- 8 M}
\begin{apts}
\item \textbf{Cost-plus pricing.} Add a fixed percentage margin to unit cost.
      Simple and safe, but ignores demand and competition. \emph{Example:}
      most retail and contract manufacturing.
\item \textbf{Value-based pricing.} Price on the customer's perceived value
      rather than on cost --- the most profitable approach where value can be
      demonstrated. \emph{Example:} Apple, whose prices reflect brand and
      ecosystem value, not the bill of materials.
\item \textbf{Penetration pricing.} Launch low to capture share quickly, then
      raise. \emph{Example:} Jio's free introductory offer, and most
      streaming services' launch pricing in India.
\item \textbf{Price skimming.} Launch high to harvest early adopters, then
      lower progressively. \emph{Example:} new smartphone and console launches.
\item \textbf{Competitive or going-rate pricing.} Set price at, above or below
      the prevailing market rate. \emph{Example:} fuel retailers and airline
      fares on competitive routes.
\item \textbf{Psychological pricing.} Exploit perception --- \Rs 999 rather than
      \Rs 1{,}000, or a high anchor price shown struck through.
      \emph{Example:} almost all e-commerce listing prices.
\item \textbf{Bundle pricing.} Sell several items together below the sum of
      their individual prices, raising average order value and clearing slower
      stock. \emph{Example:} fast-food combo meals, software suites.
\item \textbf{Freemium and subscription pricing.} Give a basic tier free and
      charge for the premium tier, or charge recurrently for access rather than
      once for ownership. \emph{Example:} Spotify, Zoom, Netflix.
\item \textbf{Dynamic pricing.} Adjust algorithmically to demand, time,
      inventory and user behaviour. \emph{Example:} Uber surge pricing, airline
      seat pricing, Amazon's continuous repricing.
\item \textbf{Promotional and loss-leader pricing.} Temporary discounts, or
      selling one item at or below cost to draw traffic that buys other items.
      \emph{Example:} festival sales; supermarket staples priced at cost.
\end{apts}
\scheme{8 M. Roughly eight strategies at 1 M each, every one requiring both an
explanation and an example. Note that pricing sits slightly outside the literal
Unit~I descriptor wording but has been examined under it, so it is worth
knowing rather than assuming it cannot be asked.}

\qq{c1b12}{A travel agency specializing in adventure tourism wants to improve
its targeted marketing. How should they create a buyer persona, and what factors
--- demographics, interests, online behaviour --- should they
consider?}{CIE-I, 3 Apr 2025; CIE-I, 16 Apr 2026 --- 10 M}
\textit{Steps to create the persona}
\begin{apts}
\item \textbf{Collect data from several sources} --- customer surveys and
      interviews, booking records, website analytics, social media insights and
      conversations with the sales team, who know the objections that never
      reach a form.
\item \textbf{Segment by shared traits} into a small number of distinct groups,
      then choose the two or three worth building personas for --- a persona per
      customer is not a persona at all.
\item \textbf{Identify goals and pain points} for each group: what they are
      trying to achieve by travelling, and what stops them booking.
\item \textbf{Build the profile} with a name, photograph, background, quote,
      goals, objections, preferred channels and buying triggers, so the whole
      team can picture one person rather than a spreadsheet row.
\item \textbf{Validate and update} against actual booking data and feedback, and
      revise as the market shifts --- a persona written once and never revisited
      becomes fiction within two seasons.
\end{apts}

\textit{Factors to consider}
\begin{apts}\setcounter{enumi}{5}
\item \textbf{Demographics.} Age roughly 20--40, gender split, income band,
      occupation --- typically salaried urban professionals and students ---
      education, city and marital status, since family stage strongly determines
      the kind of trip that is even possible.
\item \textbf{Psychographics and interests.} Trekking, scuba diving, skydiving,
      camping and cycling; fitness and outdoor lifestyle; attitude to risk;
      preference for authentic experience over luxury; environmental
      consciousness; and the values that make adventure travel attractive in the
      first place.
\item \textbf{Online behaviour.} Active on Instagram, YouTube and travel blogs;
      searches such as ``best trekking places in India'' and ``budget adventure
      trips''; consumes video and reviews before deciding; follows travel
      influencers; researches on mobile and often books on mobile too.
\item \textbf{Travel behaviour.} Frequency of travel, budget against premium
      preference, solo against group, domestic against international, planning
      lead time, and preferred season --- all of which are already recorded in
      the agency's own booking data.
\item \textbf{Pain points and objections.} Safety and equipment concerns, doubts
      about guide credentials and insurance, budget constraints, lack of
      reliable information, fitness anxiety, and distrust of agencies after
      hidden costs. These matter most, because the marketing that converts is
      the marketing that answers them.
\end{apts}

\textit{Example persona.} \textbf{``Adventure Arun''}, 28, IT professional in
Bengaluru, unmarried, income \Rs 12 lakh a year. Interests: trekking,
camping, photography. Follows travel creators on Instagram and YouTube; plans
two long weekends and one major trek a year; books online, on his phone,
usually four to six weeks ahead. Goal: unique, physically demanding experiences
he can share. Pain points: safety and guide credibility, and uncertainty about
what the quoted price actually includes.

\textit{How the persona is used.} It sets ad targeting parameters, chooses
Instagram and YouTube over print, dictates that safety credentials and
all-inclusive pricing appear prominently on landing pages, and determines that
content should be video-led rather than text-led.
\scheme{10 M. Roughly 4 M for the steps to create a persona, 4--5 M for the
factors grouped as demographics, psychographics, online behaviour, travel
behaviour and pain points, and 1--2 M for a worked example persona. The example
is what demonstrates the method has been understood rather than memorised.}

\end{qlist}
